%======================================================================
\section{Application to the Komori--Kashima Problem}\label{sec:kk}

In this section we prove Corollaries~\ref{cor:LC} and~\ref{cor:iff}.
We first observe that \kl{order} $3$ is not enough.

\subsection{Formulas of order at most \texorpdfstring{$3$}{3}}

\begin{proposition}\label{prop:order3}
Every formula of $\CL$ of \kl{order} at most $3$ belongs to $\IL$.
\end{proposition}
\begin{proof}
A formula of \kl{order} at most $2$ has the form $p_1\to\dots\to p_k\to q$ with
variables $p_1,\dots,p_k,q$, and we read it as the rule ``from
$p_1,\dots,p_k$ infer $q$''. A formula of \kl{order} at most $3$ has the form
$F=H_1\to\dots\to H_n\to r$, where the $H_i$ are of \kl{order} at most $2$ and
$r$ is a variable. Let $V$ be the closure of the empty set under the rules
$H_1,\dots,H_n$. The valuation that makes exactly the variables in $V$ true
satisfies every $H_i$, so $r\in V$ because $F\in\CL$. Following the
construction of $V$ step by step, we can derive $r$ from $H_1,\dots,H_n$ by
modus ponens in $\IL$. Hence $F\in\IL$.
\end{proof}

Hence a classical tautology of \kl{order} at most $3$ never axiomatizes a
proper extension of $\IL$.

\subsection{An axiom of Gödel--Dummett logic of order \texorpdfstring{$4$}{4}}


\begin{proof}[Proof of Corollary~\ref{cor:LC}]
The formula~\eqref{eq:SGprime} is obtained directly from
Definition~\ref{def:S}. By \cite[Lemma~4.2]{NM21}, $G'$ is \kl{minimal} in $\CL$ and $\IL\lplus G'=\LC$.
By Corollary~\ref{cor:many-logics} applied to $\IL$ and to $\CL$, we have
$\IL\lplus\St(G')=\IL\lplus G'=\LC$, and $\St(G')$ is \kl{minimal} in
$\CL$. We have $\Ord(\St(G'))\le4$ by Lemma~\ref{lem:order}, and
$\Ord(\St(G'))\ge4$ by Proposition~\ref{prop:order3}, since
$\St(G')\in\CL\setminus\IL$.
\end{proof}
\plot{The Lean file checks that the displayed formula equals
\texttt{statman Gprime 8} and has order $4$. The formula~\eqref{eq:SGprime} is generated by \texttt{paper/gen\_SGprime.py}
(supplementary material).}

\subsection{A minimality-preserving normal form}

Corollary~\ref{cor:iff} is the ``only if'' direction of the
following corollary; the ``if'' direction is trivial.

\begin{corollary}\label{cor:normal-form}
Let $L_0$ be a \kl{logic} containing $\BCI$. If $L=L_0\lplus\Gamma$, where
every $\alpha\in\Gamma$ is \kl{minimal} in $\CL$, then
$L=L_0\lplus\{\St(\alpha)\mid\alpha\in\Gamma\}$, and every $\St(\alpha)$ is
\kl{minimal} in $\CL$ and of \kl{order} at most $4$.
\end{corollary}
\begin{proof}
The first two claims follow from Corollary~\ref{cor:many-logics} applied to
$L_0$ and to $\CL$, and the last from Lemma~\ref{lem:order}.
\end{proof}

By Proposition~\ref{prop:order3} and Corollary~\ref{cor:LC}, the bound $4$
cannot be improved. Over $\IL$, Corollary~\ref{cor:normal-form} can be
sharpened slightly: since the axioms of \kl{order} at most $3$ may be
omitted by Proposition~\ref{prop:order3}, a \kl{logic} axiomatizable over
$\IL$ by $\CL$-minimal formulas is either $\IL$ or axiomatizable over $\IL$
by $\CL$-minimal formulas of \kl{order} exactly $4$.

\subsection{Antecedents of order \texorpdfstring{$3$}{3}}\label{sec:order3}

The formula $\St(G')$ is rather long. We show that a counterexample of
order $4$ must have at least two \kl{antecedents} of order $3$, which
partly explains this. For a formula $F=A_1\to\dots\to A_n\to r$ with $r$ a variable, we call
$A_1,\dots,A_n$ the \intro{antecedents} of $F$.

\begin{proposition}\label{prop:one-premise}
Let $F=A_1\to\dots\to A_n\to r\in\CL$, where $r$ is a variable and
$\Ord(F)\le4$. If at most one $A_i$ has order $3$, then either $F\in\IL$
or $\IL\lplus F=\CL$.
\end{proposition}

Note that \kl{minimality} of $F$ is not assumed.

\begin{proof}
As in the proof of Proposition~\ref{prop:order3}, we call a formula
$p_1\to\dots\to p_k\to q$ with variables $p_1,\dots,p_k,q$ ($k\ge0$) a
\intro{Horn clause} and read it as the rule ``from $p_1,\dots,p_k$ infer
$q$''; these are exactly the formulas of \kl{order} at most $2$. Since
$\Ord(F)\le4$, every $A_i$ has \kl{order} at most $3$. Recall that the \kl{antecedents} of a
formula can be permuted in $\IL$.

Let $\Gamma$ be the set of \kl{antecedents} of \kl{order} at most $2$. For a
set $X$ of variables, let $\mathrm{Cl}(X)$ be the closure of $X$ under the rules in
$\Gamma$. We identify a set of variables with the valuation that makes
exactly those variables true. Then:
\begin{enumerate}
\item[(H1)] $\mathrm{Cl}(X)$ satisfies $\Gamma$;
\item[(H2)] if $q\in \mathrm{Cl}(X)$, then $\Gamma,X\vdash_{\IL}q$, by modus ponens
  along the construction of $\mathrm{Cl}(X)$. In particular, if
  $q\in \mathrm{Cl}(\{p_1,\dots,p_k\})$, then $\Gamma\vdash_{\IL} p_1\to\dots\to p_k\to q$.
\end{enumerate}

\proofcase{Case 1: no \kl{antecedent} has \kl{order} $3$.}
Then $F$ is
$\Gamma\to r$. By (H1) and $F\in\CL$, we have $r\in \mathrm{Cl}(\emptyset)$, and hence
$F\in\IL$ by (H2).

\proofcase{Case 2: exactly one \kl{antecedent}, say
$A=B_1\to\dots\to B_m\to s$, has \kl{order} $3$.}
 Then each $B_j$ has \kl{order} at
most $2$. Assume that $F\notin\IL$, and put $V_0=\mathrm{Cl}(\emptyset)$. We first
note two facts.
\begin{itemize}
\item $r\notin V_0$, since otherwise $\Gamma\vdash_{\IL} r$ by (H2).
\item Some $B_j$ is not derivable from $\Gamma$ in $\IL$. Otherwise
  $\Gamma,A\vdash_{\IL} s$. Moreover, the valuation $\mathrm{Cl}(\{s\})$ satisfies
  $\Gamma$ and $A$, so $r\in \mathrm{Cl}(\{s\})$ because $F\in\CL$. Then
  $\Gamma,s\vdash_{\IL} r$ by (H2), and therefore $F\in\IL$, a contradiction.
\end{itemize}
Fix such a $B_j=p_1\to\dots\to p_k\to q$, and put
$V_1=\mathrm{Cl}(\{p_1,\dots,p_k\})$. Then $V_0\subseteq V_1$, and $q\notin V_1$ by
(H2). The valuation $V_1$ falsifies $B_j$ and hence satisfies $A$; it also
satisfies $\Gamma$ by (H1). Since $F\in\CL$, we obtain $r\in V_1$.

Let $H_3=\{0,\tfrac12,1\}$ be the three-element Gödel algebra, in which
$x\Rightarrow y=1$ if $x\le y$ and $x\Rightarrow y=y$ otherwise. Define a
valuation $v$ by $v(t)=1$ if $t\in V_0$, $v(t)=\tfrac12$ if
$t\in V_1\setminus V_0$, and $v(t)=0$ otherwise. A \kl{Horn clause}
$p_1\to\dots\to p_k\to q$ takes the value $1$ under $v$ if and only if
$\min_i v(p_i)\le v(q)$, where the minimum of the empty set is $1$. Now:
\begin{itemize}
\item every clause in $\Gamma$ takes the value $1$, since $V_0$ and $V_1$
  are closed under the rules in $\Gamma$;
\item $B_j$ takes the value $0$, since $v(p_i)\ge\tfrac12$ for all $i$ and
  $v(q)=0$; hence $A$ takes the value $1$;
\item $v(r)=\tfrac12$, since $r\in V_1\setminus V_0$.
\end{itemize}
Therefore $F$ takes the value $\tfrac12$, and $H_3\not\models F$. By
Jankov's criterion \cite[Theorem~1]{Jan68} (see also
\cite[Sect.~3.2]{NM21}), every implicational formula
$F\in\CL\setminus\IL$ with $H_3\not\models F$ satisfies $\IL\lplus F=\CL$.
\end{proof}
\plot{Check that \cite{Jan68} or \cite[Sect.~3.2]{NM21} covers the
implicational fragment.}

\begin{remark}\label{rem:many}
In $\St(\alpha)$, the formula $\dprem{u}$ has order $3$ if $u$ is
\kl{positive}, since $\dprem{u}=\body{u}\to\ev{u}$, and order $2$ if $u$ is
\kl{negative}, since $\dprem{u}=\ev{u}\to \body{u}$. Hence the
\kl{antecedents} of order $3$ in $\St(\alpha)$ are exactly the formulas
$\dprem{u}$ with $u$ a \kl{positive} \kl{non-leaf occurrence} of $\alpha$. By
Proposition~\ref{prop:one-premise}, if $\alpha\in\CL$ has at most one
\kl{positive} implication occurrence, then $\IL\lplus \alpha=\IL\lplus \St(\alpha)$ is
either $\IL$ or $\CL$. The formula $\St(G')$ in~\eqref{eq:SGprime} has $11$ \kl{antecedents}
of \kl{order} $3$, namely the \kl{antecedents} of the form
$(\beta\to\beta')\to\evfp{\gamma}$, where
$\gamma\in\{a\to b,\ b'\to a',\ (a'\to b')\to c'',\ (b\to a)\to c''',\
G'_7,\dots,G'_2,\ G'\}$,
whereas by Proposition~\ref{prop:one-premise} a $\CL$-minimal axiom of
\kl{order} $4$ of a proper intermediate logic must have at least two
\kl{antecedents} of \kl{order} $3$. We do not know the least
possible number.
\end{remark}
